
\section{Quantitatively calibrated seven-ray completion}
\label{sec:xxvii-calibrated}

Edition XXVI proves pressure-quotient realizability using large
source-independent gauge constants. Here we substantially reduce
those constants and prove a sharp zero-source floor. All statements
are exact on a periodic-axial cylinder, or on an embedded annular
disk of a three-torus; no assertion of a compact axial gauge on
$\mathbb R^3$ is made.

\begin{theorem}[Improved universal gauge with an arbitrary cone floor]
\label{thm:xxvii-universal}
Let $Q(r,z)$ be a smooth symmetric co-rotating tensor, radially
compactly supported away from the axis, and let $\chi(r)$ be a
smooth radial cutoff equal to one throughout the radial
projection of $\operatorname{supp}Q$. Write
\[
M=\max_{i,j}\|Q_{ij}\|_{L^\infty},\qquad
\mu>0,\qquad
a=\frac{3}{8}\mu+\frac{13}{4}M,\qquad
b=\frac{85}{8}\mu+\frac{471}{4}M,\qquad
h=\mu+8M.
\]
Define the seven coefficients (ray aperture $1/4$)
\begin{align*}
d_1&=Q_{rr}-16(Q_{\theta\theta}+Q_{zz})
          +(\mu+33M)\chi,\\
d_2&=8Q_{\theta\theta}+2Q_{r\theta}
          +(\mu+10M)\chi,\\
d_3&=8Q_{\theta\theta}-2Q_{r\theta}
          +(3\mu+26M)\chi,\\
d_4&=8Q_{zz}+2Q_{rz}-8Q_{\theta z}
          +(2\mu+18M)\chi,\\
d_5&=8Q_{zz}-2Q_{rz}+8Q_{\theta z}
          +(2\mu+18M)\chi,\\
d_6&=8Q_{\theta z}+(\mu+8M)\chi,\\
d_7&=-8Q_{\theta z}+(\mu+8M)\chi.
\end{align*}
Then, everywhere,
\begin{equation}\label{eq:xxvii-completion}
\boxed{\quad
d_j\ge\mu\chi,\qquad
\sum_{j=1}^7d_jv_j\otimes v_j
=Q+a\chi I+b\chi e_r\otimes e_r.\quad}
\end{equation}
The carrier is globally pressure-neutral in the same periodic
geometries as Edition XXVI:
\[
\nabla\cdot\bigl(\sum_jd_jv_j\otimes v_j-Q\bigr)
=\nabla\left((a+b)\chi+
  b\int_0^r\frac{\chi(s)}s\,ds\right).
\]
In particular, every matrix entry of the constructed stress obeys
\begin{equation}\label{eq:xxvii-size}
\boxed{\quad
\left\|\sum_jd_jv_j\otimes v_j\right\|_{\ell^\infty_{\rm entry}}
\le122M+11\mu.\quad}
\end{equation}
\end{theorem}
\begin{proof}
The seven ray tensor identities are those in
Section~\ref{sec:xxvi-global-ray}. For a general
seven-coefficient inverse with middle parameter $h\chi$,
the off-diagonal and diagonal equations give
\[
\begin{aligned}
d_6&=h\chi+8Q_{\theta z},&d_7&=h\chi-8Q_{\theta z},\\
d_2&=8Q_{\theta\theta}+2Q_{r\theta}
      +(8a-2h)\chi,&
d_3&=8Q_{\theta\theta}-2Q_{r\theta}+8a\chi,\\
d_4&=8Q_{zz}+2Q_{rz}-8Q_{\theta z}
      +(8a-h)\chi,&
d_5&=8Q_{zz}-2Q_{rz}+8Q_{\theta z}
      +(8a-h)\chi,\\
d_1&=Q_{rr}-16(Q_{\theta\theta}+Q_{zz})
      +(b-31a+2h)\chi.
\end{aligned}
\]
The chosen values satisfy exactly
\[
8a-2h=\mu+10M,\quad 8a=3\mu+26M,\quad
8a-h=2\mu+18M,\quad
b-31a+2h=\mu+33M.
\]
Where $Q\ne0$, $\chi=1$ and $|Q_{ij}|\le M$.
The worst-case signed contributions to $d_1,\ldots,d_7$
are respectively $-33M,-10M,-10M,-18M,-18M,-8M,-8M$.
Every coefficient therefore exceeds $\mu$. In a cutoff
collar $Q=0$, the displayed positive constants multiply $\chi$,
so $d_j\ge\mu\chi$ globally. The tensor identity follows by
substitution in the six synthesis equations.

The carrier divergence is a gradient as in Edition XXVI.
The only potentially largest entry of
$S=Q+a\chi I+b\chi e_r\otimes e_r$ is the radial diagonal:
\[
|S_{rr}|\le M+a+b=M+121M+11\mu=122M+11\mu.
\]
The other diagonals are at most
$(17/4)M+(3/8)\mu$, and the off-diagonals at most $M$.
\end{proof}

\begin{proposition}[Sharp zero-source cost, and minimax within a fixed-splitting inverse]
\label{prop:xxvii-sharp-floor}
For $Q=0$, $\chi=1$, any strict seven-ray representation
\[
aI+b e_r\otimes e_r=\sum_jd_jv_j\otimes v_j,\qquad
d_j\ge\mu>0
\]
requires
\begin{equation}\label{eq:xxvii-sharp-floor}
\boxed{\quad b\ge\frac{85}{8}\mu.\quad}
\end{equation}
Equality is attained by
$(d_1,\ldots,d_7)=(\mu,\mu,3\mu,2\mu,2\mu,\mu,\mu)$.

Moreover, within the source-independent fixed-$h$ inverse
displayed in the proof of Theorem~\ref{thm:xxvii-universal},
with fixed scalar $a,b,h$ and arbitrary symmetric
$|Q_{ij}|\le M$, one necessarily has
\[
h\ge\mu+8M,\qquad
a\ge\frac{\mu+2h+10M}{8},\qquad
b\ge\mu+31a-2h+33M,
\]
which imply
\begin{equation}\label{eq:xxvii-fixed-minimax}
\boxed{\quad
b\ge\frac{85}{8}\mu+\frac{471}{4}M.\quad}
\end{equation}
The construction of Theorem~\ref{thm:xxvii-universal}
attains this bound within the stated fixed-inverse class.
This is \emph{not} a claim that the pressure quotient has
the same optimal constant among all nonlinear
source-dependent coefficient selections.
\end{proposition}
\begin{proof}
At $Q=0$, the three off-diagonal tensor constraints give
$d_6=d_7=h$, $d_3=d_2+2h$, $d_4=d_5$.
Equality of the two transverse diagonal entries gives
$d_4=d_5=d_2+h$. Hence
\[
a=(d_2+2h)/8,\qquad
b=d_1+\frac{31}{8}d_2+\frac{23}{4}h.
\]
Since $d_1,d_2,h\ge\mu$, this gives the sharp $85\mu/8$
bound and its equality case.

For the fixed inverse in the theorem, testing all signed
corners of the independent box $|Q_{ij}|\le M$ gives
the three inequalities displayed, as well as
$8a\ge\mu+10M$ and $8a-h\ge\mu+18M$.
The first three dominate the latter two for
$h\ge\mu+8M$. Eliminating $a$ and then $h$ yields
\[
b\ge\frac{39}{8}\mu+\frac{287}{4}M
                 +\frac{23}{4}h
\ge\frac{85}{8}\mu+\frac{471}{4}M.
\]
The displayed explicit coefficients attain equality.
\end{proof}

\paragraph{Independent audit, baseline, and real source limitation.}
A distinct symbolic reconstruction checks every one of the
six symmetric tensor components, and $20{,}000$ exact-rational
box tests verify all seven coefficient floors.
An independent linear-programming calculation finds the
fixed-inverse optimum $(a,h,b)=(29/8,9,1027/8)$ at
$(\mu,M)=(1,1)$, and independently confirms the
sharp zero-source constant $85/8$.
The $M$ slope of the carrier decreases from the source-independent
Edition XXVI value $528$ to $471/4=117.75$.
The original A.11/A.19 source remains unauthenticated;
no H8 smallness threshold or all-mode Navier--Stokes
closure follows from a universal pressure-quotient identity.
The true end-to-end unforced theorem and Lean PDE proof remain open.
